%% ============================================================
%%  Nota Técnica de Metodologia — SEAFLOR 2026
%%  HHI · Decomposição Comercial · Estimativa de Receita Foregone
%%  Observatório FIESC | Junho 2026
%%  Compile em: overleaf.com (pdflatex)
%% ============================================================
\documentclass[12pt, a4paper]{article}

%% ── Pacotes ──────────────────────────────────────────────────
\usepackage[utf8]{inputenc}
\usepackage[T1]{fontenc}
\usepackage[brazil]{babel}
\usepackage{geometry}
\usepackage{amsmath, amssymb, mathtools}
\usepackage{booktabs}
\usepackage{array}
\usepackage{multirow}
\usepackage{xcolor}
\usepackage{graphicx}
\usepackage{enumitem}
\usepackage{lmodern}
\usepackage{caption}
\usepackage{float}
\usepackage{siunitx}
\usepackage{makecell}
\usepackage[most]{tcolorbox}
\usepackage{tikz}
\usetikzlibrary{shapes.geometric, arrows.meta, positioning}
\usepackage{titlesec}
\usepackage[colorlinks=true, linkcolor=azulFIESC, urlcolor=azulMed,
            citecolor=azulFIESC, pdftitle={Nota Técnica SEAFLOR 2026}]{hyperref}

\geometry{left=3cm, right=2.5cm, top=3cm, bottom=2.5cm}

%% ── Cores ────────────────────────────────────────────────────
\definecolor{azulFIESC}{HTML}{1F4E79}
\definecolor{azulMed}{HTML}{2E75B6}
\definecolor{azulCl}{HTML}{DEEAF1}
\definecolor{verde}{HTML}{375623}
\definecolor{vermelho}{HTML}{C00000}
\definecolor{cinzaTexto}{HTML}{595959}
\definecolor{laranjaEsc}{HTML}{C55A11}
\definecolor{amareloCl}{HTML}{FFFDE7}
\definecolor{verdeCl}{HTML}{E8F5E9}

%% ── Caixas tcolorbox ─────────────────────────────────────────
\tcbuselibrary{most}

\newtcolorbox{formulabox}[1][]{%
  colback=azulCl, colframe=azulFIESC,
  fonttitle=\bfseries\color{white}, coltitle=white,
  boxrule=1.5pt, arc=4pt, title={#1}, left=6pt, right=6pt,
  before upper={\setlength{\abovedisplayskip}{4pt}%
                \setlength{\belowdisplayskip}{4pt}}}

\newtcolorbox{exemplobox}[1][Exemplo numérico]{%
  colback=amareloCl, colframe=laranjaEsc,
  fonttitle=\bfseries\color{laranjaEsc},
  title={#1}, arc=4pt, boxrule=1pt}

\newtcolorbox{atencaobox}[1][]{%
  colback=vermelho!5, colframe=vermelho,
  fonttitle=\bfseries\color{vermelho},
  title={\ifx&#1&Atenção\else#1\fi}, arc=4pt, boxrule=1pt}

\newtcolorbox{interpretabox}[1][]{%
  colback=verdeCl, colframe=verde,
  fonttitle=\bfseries\color{verde},
  title={\ifx&#1&Como interpretar\else#1\fi}, arc=4pt, boxrule=1pt}

%% ── Títulos de seção ─────────────────────────────────────────
\titleformat{\section}
  {\large\bfseries\color{azulFIESC}}{\thesection}{1em}{}[\color{azulMed}\hrule\vspace{2pt}]
\titleformat{\subsection}
  {\normalsize\bfseries\color{azulMed}}{\thesubsection}{1em}{}
\titleformat{\subsubsection}
  {\normalsize\itshape\color{cinzaTexto}}{\thesubsubsection}{1em}{}

%% ── siunitx ──────────────────────────────────────────────────
\sisetup{group-separator={.}, group-minimum-digits=4,
         output-decimal-marker={,}}

%% ============================================================
\begin{document}

%% ── CAPA ─────────────────────────────────────────────────────
\begin{titlepage}
  \centering
  \vspace*{1.5cm}
  {\color{azulFIESC}\rule{\linewidth}{2.5pt}}\\[0.5cm]
  {\LARGE\bfseries\color{azulFIESC}
    Nota Técnica de Metodologia\\[0.4cm]
    SEAFLOR 2026\\[0.3cm]
    \large Concentração de Mercado (HHI),\\
    Decomposição Comercial e\\
    Estimativa de Receita Foregone
  }\\[0.5cm]
  {\color{azulFIESC}\rule{\linewidth}{1pt}}\\[1.5cm]
  {\large\color{cinzaTexto}
    Observatório FIESC --- Competitividade Industrial\\[0.3cm]
    Junho de 2026
  }\\[2.5cm]
  \begin{minipage}{0.9\textwidth}
    \centering\itshape\color{cinzaTexto}
    Este documento apresenta, de forma didática e com rigor matemático,
    as três metodologias centrais do relatório SEAFLOR~2026.
    Para cada metodologia são descritos: a lógica conceitual,
    as equações com todos os termos definidos,
    as premissas e limitações, e um exemplo numérico
    passo a passo com dados reais do Complexo Florestal de SC.\\[0.8cm]
    \normalfont\small
    \textbf{Metodologia 1 ---} Decomposição do impacto tarifário:
      destruição bruta, desvio de comércio, destruição líquida e efeito-preço.\\[0.3cm]
    \textbf{Metodologia 2 ---} Concentração geográfica de exportações:
      índice Herfindahl-Hirschman (HHI).\\[0.3cm]
    \textbf{Metodologia 3 ---} Estimativa de receita não realizada:
      análise contrafactual por cenários.
  \end{minipage}\\[1.5cm]
  {\footnotesize\color{cinzaTexto}
    Fonte dos dados: ComexStat/MDIC\quad$\cdot$\quad
    Classificação: NCM--CNAE FIESC\quad$\cdot$\quad
    Período: Jan--Mai 2025 vs.\ Jan--Mai 2026
  }
\end{titlepage}

\tableofcontents
\newpage

%%=============================================================
\section{Apresentação e Escopo}

Este documento técnico descreve as metodologias quantitativas utilizadas
pelo Observatório FIESC para mensurar o impacto das tarifas americanas
de~2025--2026 sobre as exportações do \textbf{Complexo Florestal de
Santa Catarina}.

O \textbf{Complexo Florestal} compreende quatro setores classificados pela CNAE:

\begin{center}
\renewcommand{\arraystretch}{1.3}
\begin{tabular}{>{\bfseries}cl}
\toprule
CNAE & Atividade econômica \\
\midrule
2  & Produção Florestal --- base florestal (toras, lenha, carvão vegetal) \\
16 & Fabricação de Produtos de Madeira (serrada, compensada, MDF, molduras) \\
17 & Fabricação de Celulose, Papel e Produtos de Papel \\
31 & Fabricação de Móveis e Componentes \\
\bottomrule
\end{tabular}
\end{center}

\noindent O documento está dividido em cinco seções principais:

\begin{description}[style=nextline, leftmargin=2cm]
  \item[Seção~\ref{sec:bases}]
    Bases de dados, fluxo de processamento e recortes analíticos.
  \item[Seção~\ref{sec:decomp}]
    Metodologia de decomposição do impacto tarifário.
  \item[Seção~\ref{sec:hhi}]
    Índice de concentração Herfindahl-Hirschman (HHI).
  \item[Seção~\ref{sec:receita}]
    Estimativa da receita foregone por análise contrafactual.
  \item[Seção~\ref{sec:resultados}]
    Resultados por CNAE com as três metodologias.
\end{description}

%%=============================================================
\section{Bases de Dados, Fluxo de Processamento e Recortes Analíticos}
\label{sec:bases}

\subsection{Fonte Primária: ComexStat/MDIC}

Todos os dados utilizados provêm do \textbf{ComexStat}, sistema do Ministério
do Desenvolvimento, Indústria, Comércio e Serviços (MDIC). O ComexStat
registra, ao nível de oito dígitos da Nomenclatura Comum do Mercosul (NCM),
cada operação declarada ao Siscomex, com as seguintes variáveis relevantes:

\begin{itemize}[noitemsep]
  \item \textbf{NCM} (8 dígitos): classificação do produto exportado;
  \item \textbf{UF} do exportador: filtro para \texttt{sg\_uf = 'SC'};
  \item \textbf{País de destino}: código ISO do país comprador;
  \item \textbf{Mês/Ano}: período da operação;
  \item \textbf{Valor FOB} (US\$): valor na fronteira de embarque;
  \item \textbf{Peso líquido} (kg): volume físico exportado.
\end{itemize}

Os arquivos são disponibilizados em dois formatos:
\begin{itemize}[noitemsep]
  \item \textbf{Histórico consolidado} (\textit{gold}): anos 2022--2025 em Parquet.
  \item \textbf{Bronze mensal}: arquivos \texttt{EXP\_AAAAMM.csv} de 2026,
    publicados mensalmente pelo MDIC.
\end{itemize}

\subsection{Dicionário de Classificação NCM--CNAE}

A correspondência entre o código NCM (produto) e a CNAE (setor) é feita por
um dicionário desenvolvido pela FIESC e atualizado em abril de 2026.
Cada NCM recebe \emph{exatamente um} código CNAE (\textit{mapeamento 1:1}):

\begin{center}
\renewcommand{\arraystretch}{1.3}
\begin{tabular}{cll}
\toprule
\textbf{CNAE} & \textbf{Capítulos SH principais} & \textbf{Produtos representativos} \\
\midrule
2  & SH 44 (parcial), SH 45 & Toras, lenha, carvão vegetal \\
16 & SH 44 (principal)       & Madeira serrada, compensada, MDF, molduras \\
17 & SH 47, 48, 49           & Celulose, papel kraft, embalagens \\
31 & SH 94 (parcial)         & Móveis, assentos, componentes \\
\bottomrule
\end{tabular}
\end{center}

\subsection{Fluxo de Processamento de Dados}

O diagrama abaixo resume a trajetória dos dados da fonte bruta até os
indicadores do relatório:

\vspace{1em}
\begin{center}
\begin{tikzpicture}[
  node distance=0.7cm and 1.6cm,
  box/.style={draw, rounded corners=4pt, fill=azulCl, text centered,
              minimum height=1cm, minimum width=2.8cm, font=\small\color{azulFIESC}},
  arr/.style={-Stealth, thick, color=azulFIESC},
  lbl/.style={font=\scriptsize\itshape, color=cinzaTexto, text centered}]

  \node[box] (src)  {ComexStat\\NCM 8-dig.};
  \node[box, right=of src]  (fil)  {Filtro SC\\+ CNAE dict.};
  \node[box, right=of fil]  (agg)  {Agrega por\\país/setor/mês};
  \node[box, right=of agg]  (calc) {Cálculo dos\\indicadores};
  \node[box, right=of calc] (out)  {Excel\\SEAFLOR};

  \draw[arr] (src)  -- node[above,lbl]{CSV bronze} (fil);
  \draw[arr] (fil)  -- node[above,lbl]{silver} (agg);
  \draw[arr] (agg)  -- node[above,lbl]{gold Parquet} (calc);
  \draw[arr] (calc) -- node[above,lbl]{Python/pandas} (out);
\end{tikzpicture}
\end{center}
\vspace{0.5em}

\begin{enumerate}
  \item \textbf{Fonte bruta (bronze):} Arquivos CSV mensais com todas as operações
    de exportação do Brasil no ComexStat.
  \item \textbf{Filtro e enriquecimento (silver):} Seleção de \texttt{sg\_uf = 'SC'}
    e aplicação do dicionário NCM--CNAE para atribuir cada NCM ao seu setor.
  \item \textbf{Agregação (gold):} Dados agregados por
    \{país de destino, CNAE, mês/ano\}, armazenados em Parquet.
  \item \textbf{Cálculo dos indicadores:} Scripts Python (DuckDB + pandas)
    calculam as métricas de desvio, HHI e receita foregone.
  \item \textbf{Saída:} Resultados escritos no Excel
    \texttt{SEAFLOR\_2026\_JanMai\_3anos\_v2.xlsx}.
\end{enumerate}

\subsection{Recortes Analíticos}

Três recortes definem o universo de análise:

\begin{description}[style=nextline]
  \item[Recorte temporal]
    Comparação \textbf{Janeiro a Maio} de cada ano --- período simétrico que elimina
    sazonalidade. O período Jan--Mai 2026 é o primeiro semestre completo sob o
    regime de tarifas de 50\% (vigentes desde 6 de agosto de 2025).

  \item[Recorte setorial]
    Complexo Florestal: CNAE~2 + 16 + 17 + 31. O ``Complexo'' consolida pelo
    filtro de CNAE; os segmentos individuais podem usar filtros adicionais
    de competitividade SC (\texttt{sc\_competitiva}), gerando pequenas divergências
    entre a soma dos setores e o total do Complexo.

  \item[Recorte geográfico]
    \textbf{EUA} = código de país 249 (Estados Unidos).
    \textbf{OUT} = todos os demais países (complemento de EUA).
    Para o HHI, o universo abrange todos os países com exportação positiva no período.
\end{description}

\subsection{Notação}

A notação utilizada ao longo do documento:

\begin{align*}
  X_{d,s}^{t}
    &= \text{valor FOB (US\$) exportado pelo setor }s
       \text{ para o destino }d\text{ em Jan--Mai do ano }t \\
  Q_{d,s}^{t}
    &= \text{volume (kg líq.) exportado pelo setor }s
       \text{ para o destino }d\text{ em Jan--Mai do ano }t \\
  X_{\text{TOT},s}^{t}
    &= \textstyle\sum_{d} X_{d,s}^{t}
       \quad\text{(total exportado pelo setor }s\text{ em Jan--Mai de }t\text{)} \\
  \text{EUA}
    &= \text{destino ``Estados Unidos''} \\
  \text{OUT}
    &= \text{todos os destinos exceto EUA}
\end{align*}

%%=============================================================
\section{Metodologia 1 --- Decomposição do Impacto Tarifário}
\label{sec:decomp}

\subsection{Lógica Conceitual}

A decomposição adapta ao contexto tarifário a distinção clássica de
\textbf{Jacob Viner} (1950) entre \textit{trade creation} e
\textit{trade diversion}, originalmente desenvolvida para uniões aduaneiras.
A lógica central pode ser descrita em quatro movimentos:

\begin{enumerate}
  \item Quando os EUA impõem tarifa, o exportador catarinense perde
    competitividade naquele mercado --- o comércio com os EUA é
    \textbf{destruído}.
  \item Parte dessa destruição pode ser compensada se o exportador
    redireciona produção a outros destinos ---
    há \textbf{desvio de comércio} para fora dos EUA.
  \item O que sobra após o desvio é a \textbf{destruição líquida}:
    a perda real e não compensada de receita em \textit{valor}.
  \item Mas mesmo quando há desvio, pode haver perda de receita
    \textit{por preço}: o mercado americano paga preços unitários
    superiores. Se a produção desviada obtém preço inferior, a receita
    cai mesmo que o volume não caia --- esse é o \textbf{efeito-preço}.
\end{enumerate}

\begin{formulabox}[Fluxo lógico da decomposição]
\[
\underbrace{DB_s}_{\substack{\text{queda nas}\\[2pt]\text{vendas aos EUA}}}
\;-\;
\underbrace{DC_s}_{\substack{\text{recuperado via}\\[2pt]\text{outros mercados}}}
\;=\;
\underbrace{DL_s}_{\substack{\text{perda não}\\[2pt]\text{compensada (valor)}}}
\;+\;
\underbrace{EP_s}_{\substack{\text{perda por preço}\\[2pt]\text{unitário inferior}}}
\;=\;
\underbrace{PT_s}_{\substack{\text{perda total}\\[2pt]\text{estimada}}}
\]
\end{formulabox}

\subsection{Destruição Bruta (\texorpdfstring{$DB_s$}{DB})}

\begin{formulabox}[Destruição Bruta]
\begin{equation}
  DB_s \;=\; X_{\text{EUA},s}^{2025} \;-\; X_{\text{EUA},s}^{2026}
  \label{eq:db}
\end{equation}
\end{formulabox}

\noindent\textbf{O que mede:} a queda absoluta em dólares nas exportações do
setor~$s$ para os EUA entre Jan--Mai 2025 e Jan--Mai 2026.

\noindent\textbf{Interpretação:}
\begin{itemize}[noitemsep]
  \item $DB_s > 0$: houve queda nas exportações para os EUA.
    \textit{Caso de Madeira (CNAE~16) e Móveis (CNAE~31).}
  \item $DB_s < 0$: as exportações para os EUA \emph{cresceram}.
    Esse caso ocorre na Base Florestal (CNAE~2), possivelmente por
    contratos de longo prazo firmados antes das tarifas.
\end{itemize}

\begin{atencaobox}
$DB_s$ mede o efeito \emph{sobre o canal EUA}, não sobre o total exportado.
Um setor pode exportar mais no total e ainda ter $DB_s > 0$, se a perda
nos EUA for compensada por outros destinos.
\end{atencaobox}

\subsection{Desvio de Comércio (\texorpdfstring{$DC_s$}{DC})}

\begin{formulabox}[Desvio de Comércio]
\begin{equation}
  DC_s \;=\; \max\!\bigl(X_{\text{OUT},s}^{2026} \;-\; X_{\text{OUT},s}^{2025},\; 0\bigr)
  \label{eq:dc}
\end{equation}
\end{formulabox}

\noindent\textbf{O que mede:} o ganho líquido nas exportações para todos os
destinos \emph{exceto} os EUA. Se o setor expandiu vendas ao resto do mundo,
esse incremento representa a parcela da produção antes destinada aos EUA
que foi redirecionada com sucesso.

\noindent\textbf{Por que $\max(\cdot,0)$?}
Se as exportações para fora dos EUA \emph{caíram}, isso é uma perda adicional
--- não uma compensação. Nesses casos, $DC_s = 0$ e toda a destruição bruta
é líquida.

\begin{interpretabox}
Um $DC_s$ elevado indica resiliência exportadora: o setor encontrou mercados
alternativos. Um $DC_s \approx 0$ indica que a perda nos EUA não foi
compensada --- o volume antes destinado a esse mercado não foi absorvido
por outros compradores.
\end{interpretabox}

\subsection{Destruição Líquida (\texorpdfstring{$DL_s$}{DL})}

\begin{formulabox}[Destruição Líquida]
\begin{equation}
  DL_s \;=\; DB_s \;-\; DC_s
  \label{eq:dl}
\end{equation}
\end{formulabox}

\noindent\textbf{O que mede:} o impacto líquido das tarifas sobre o
\emph{valor total} exportado, após descontar a compensação obtida em
outros mercados.

\begin{itemize}[noitemsep]
  \item $DL_s > 0$: a perda nos EUA supera o ganho em outros mercados
    --- caso de Madeira e Móveis.
  \item $DL_s < 0$: os ganhos em outros mercados superaram a perda nos EUA
    --- o setor \emph{cresceu em valor total}. Caso do Papel e Celulose
    (CNAE~17), que expandiu exportações para Argentina e Europa.
\end{itemize}

\subsection{Prêmio de Mercado EUA e Efeito-Preço}
\label{sec:ep}

Mesmo quando o desvio compensa a perda em volume, o setor pode estar
exportando a um \emph{preço unitário inferior} ao praticado com os EUA.
Esse diferencial é o \textbf{Prêmio de Mercado EUA} ($\pi_s$).

O valor unitário médio de um destino $d$ é:

\begin{equation}
  VU_{d,s}^{t} \;=\; \frac{X_{d,s}^{t}}{Q_{d,s}^{t}}
  \quad\text{(US\$/kg)}
  \label{eq:vu}
\end{equation}

\begin{formulabox}[Prêmio de Mercado EUA e Efeito-Preço]
\begin{align}
  \pi_s
    &\;=\; VU_{\text{EUA},s}^{2025} \;-\; VU_{\text{OUT},s}^{2025}
    \label{eq:premio} \\[6pt]
  EP_s
    &\;=\; \max\!\bigl(\Delta Q_{\text{OUT},s},\;0\bigr)
           \;\times\;
           \max\!\bigl(\pi_s,\;0\bigr)
    \label{eq:ep}
\end{align}
onde $\;\Delta Q_{\text{OUT},s} = Q_{\text{OUT},s}^{2026} - Q_{\text{OUT},s}^{2025}$.
\end{formulabox}

\noindent\textbf{Intuição:} se o setor enviou mais quilos para outros mercados
($\Delta Q > 0$), mas a um preço inferior ao americano ($\pi_s > 0$), houve
perda de receita por unidade desviada. O Efeito-Preço quantifica esse
custo de oportunidade.

\noindent\textbf{Por que duas condições $\max(\cdot,0)$?}
\begin{enumerate}[noitemsep]
  \item Só conta o volume \emph{adicional} efetivamente desviado
    (se não houve aumento de volume, não há perda de preço unitário).
  \item Só prêmios positivos geram efeito
    (se o mercado alternativo pagasse mais que os EUA, não haveria custo de oportunidade).
\end{enumerate}

\begin{atencaobox}
O Efeito-Preço usa o prêmio de \textbf{2025} como referência, não de~2026.
Isso evita contaminação pelo próprio choque (as tarifas podem ter alterado
os preços dos EUA em 2026). O prêmio de~2025 representa o equilíbrio
\emph{pré-choque}.
\end{atencaobox}

\subsection{Perda Total (\texorpdfstring{$PT_s$}{PT})}

\begin{formulabox}[Perda Total --- Método A]
\begin{equation}
  PT_s
    \;=\; DL_s + EP_s
    \;=\; (DB_s - DC_s)
          + \max(\Delta Q_{\text{OUT},s},\,0)\cdot\max(\pi_s,\,0)
  \label{eq:pt}
\end{equation}
\end{formulabox}

A Perda Total combina:
\begin{enumerate}[noitemsep]
  \item A perda líquida de \emph{valor} exportado ($DL_s$):
    quanto o setor deixou de vender no total.
  \item A perda de receita por \emph{diferencial de preço} ($EP_s$):
    quanto a mais o setor teria faturado se o volume desviado
    fosse vendido ao preço americano.
\end{enumerate}

\subsection{Exemplo Passo a Passo: Madeira (CNAE~16)}

\begin{exemplobox}[Exemplo passo a passo --- Madeira (CNAE 16)]
\small
\textbf{Dados observados (Jan--Mai 2025 vs.\ 2026):}
\begin{align*}
  X_{\text{EUA},16}^{2025} &= \text{US\$~262,1~mi} &
  X_{\text{EUA},16}^{2026} &= \text{US\$~143,9~mi} \\
  X_{\text{OUT},16}^{2025} &= \text{US\$~277,4~mi} &
  X_{\text{OUT},16}^{2026} &= \text{US\$~309,9~mi} \\
  VU_{\text{EUA},16}^{2025} &= 0{,}942 \text{ US\$/kg} &
  VU_{\text{OUT},16}^{2025} &= 0{,}434 \text{ US\$/kg} \\
  Q_{\text{OUT},16}^{2025} &= 638{,}4\text{ mil t} &
  Q_{\text{OUT},16}^{2026} &= 712{,}0\text{ mil t}
\end{align*}

\medskip
\textbf{Passo 1 --- Destruição Bruta:}
\[
  DB_{16} = 262{,}1 - 143{,}9 = \mathbf{118{,}1\ \text{mi}}
\]

\textbf{Passo 2 --- Desvio de Comércio:}
\[
  DC_{16} = \max(309{,}9 - 277{,}4,\;0) = \max(32{,}5,\;0) = \mathbf{32{,}5\ \text{mi}}
\]

\textbf{Passo 3 --- Destruição Líquida:}
\[
  DL_{16} = 118{,}1 - 32{,}5 = \mathbf{85{,}6\ \text{mi}}
\]

\textbf{Passo 4 --- Prêmio de Mercado EUA:}
\[
  \pi_{16} = 0{,}942 - 0{,}434 = \mathbf{0{,}508\ \text{US\$/kg}}
  \quad(+116{,}8\%\text{ em relação ao preço nos demais mercados})
\]

\textbf{Passo 5 --- Efeito-Preço:}
\[
  \Delta Q_{16} = 712{,}0 - 638{,}4 = +73{,}6\text{ mil ton}
  \quad\Rightarrow\quad
  EP_{16} \approx 73{.}600 \times 0{,}508 \approx \mathbf{30{,}1\ \text{mi}}
\]
\emph{(O valor exato de US\$~30,1~mi resulta do cálculo por NCM, que captura
a variação de mix de produto dentro da Madeira.)}

\textbf{Passo 6 --- Perda Total:}
\[
  PT_{16} = 85{,}6 + 30{,}1 = \mathbf{US\$\ 115{,}7\ \text{milhões}}
\]
\end{exemplobox}

%%=============================================================
\section{Metodologia 2 --- Concentração de Mercado: Índice HHI}
\label{sec:hhi}

\subsection{Origem e Intuição}

O \textbf{Índice Herfindahl-Hirschman (HHI)} foi proposto independentemente
por Orris Herfindahl (1950) e Albert Hirschman (1945) para medir concentração
em estruturas de mercado. No comércio internacional, é amplamente usado para
quantificar a \textbf{diversificação geográfica das exportações}:
quanto mais distribuídas entre países, menor o HHI e menor a vulnerabilidade
a choques em um único mercado.

\medskip
A \textbf{ideia central} do HHI é simples:
o índice soma o quadrado das participações percentuais de cada destino.
Elevar ao quadrado penaliza desproporcionalmente destinos com participação
elevada --- um país com 40\% do total contribui com $40^2 = 1.600$ pontos,
enquanto quatro países com 10\% cada contribuem juntos com apenas
$4 \times 10^2 = 400$ pontos.

Isso significa que \textbf{reduzir a participação de um destino dominante}
é muito mais eficaz para baixar o HHI do que adicionar muitos destinos pequenos.

\subsection{Formulação Matemática}

\begin{formulabox}[Definição --- Índice HHI]
\begin{equation}
  HHI_{s,t}
    \;=\; \sum_{i=1}^{N_{s,t}}
          \left(\frac{X_{i,s}^{t}}{X_{\text{TOT},s}^{t}} \times 100\right)^{\!2}
  \label{eq:hhi}
\end{equation}

\begin{itemize}[noitemsep]
  \item $i$: índice do país de destino ($i = 1, \dots, N_{s,t}$)
  \item $X_{i,s}^{t}$: valor FOB exportado para o país $i$, setor $s$,
    em Jan--Mai do ano $t$
  \item $X_{\text{TOT},s}^{t}$: valor total exportado pelo setor $s$
    em Jan--Mai do ano $t$
  \item $N_{s,t}$: número de países com exportação positiva
\end{itemize}
\end{formulabox}

O índice varia de 0 (distribuição perfeita entre infinitos destinos)
a 10.000 (exportação para um único país).

\subsection{Classificação de Concentração}

A classificação adotada segue os critérios do
Departamento de Justiça dos EUA (DOJ) e da Federal Trade Commission (FTC),
adaptados ao contexto de exportações:

\begin{center}
\renewcommand{\arraystretch}{1.3}
\begin{tabular}{cll}
\toprule
\textbf{HHI} & \textbf{Classificação} & \textbf{Interpretação para exportações} \\
\midrule
$\geq 2.500$        & Alta concentração      & Alta dependência; risco elevado de choque \\
$1.500$ -- $2.499$  & Concentração moderada  & Vulnerabilidade relevante a destino principal \\
$1.000$ -- $1.499$  & Baixa concentração     & Pauta diversificada com alguma exposição \\
$< 1.000$           & Diversificado          & Alta dispersão; baixa vulnerabilidade \\
\bottomrule
\end{tabular}
\end{center}

\subsection{Contribuição Individual ao HHI}

Para identificar quais países mais \emph{pressionam} o índice,
calcula-se a contribuição individual do país $i$:

\begin{equation}
  h_{i,s}^{t}
    \;=\; \left(\frac{X_{i,s}^{t}}{X_{\text{TOT},s}^{t}} \times 100\right)^{\!2}
  \label{eq:contrib}
\end{equation}

tal que $HHI_{s,t} = \displaystyle\sum_{i=1}^{N} h_{i,s}^{t}$.

Essa decomposição permite calcular exatamente quanto o HHI mudaria
se a participação de um determinado destino fosse alterada,
sem necessidade de recalcular o índice completo.

\subsection{Exemplo Passo a Passo: Complexo Florestal SC}

\begin{exemplobox}[Cálculo do HHI --- Complexo Florestal SC, Jan-Mai 2025 vs.\ 2026]
\small
\textbf{Jan--Mai 2025} (total exportado: US\$~822,2~mi):

\medskip
\begin{center}
\renewcommand{\arraystretch}{1.25}
\begin{tabular}{lS[table-format=3.1]S[table-format=2.2]S[table-format=4.1]}
\toprule
\textbf{Destino} & \textbf{Valor (US\$ mi)} & \textbf{Share (\%)} & \textbf{Share\textsuperscript{2}} \\
\midrule
Estados Unidos &  322.2 & 39.18 & 1535.1 \\
Argentina      &   46.8 &  5.69 &   32.4 \\
México         &   38.3 &  4.66 &   21.7 \\
Reino Unido    &   25.1 &  3.05 &    9.3 \\
Alemanha       &   20.5 &  2.49 &    6.2 \\
Outros         &  369.3 & 44.93 & \textit{(vários)} \\
\midrule
\textbf{HHI total} & & & \textbf{1.298} \\
\bottomrule
\end{tabular}
\end{center}

\medskip
\noindent Apenas os EUA contribuem com \textbf{1.535 pontos} de um total de
1.298 do HHI. Isso é possível porque os demais países, individualmente
pequenos, contribuem positivamente mas com valores baixos, resultando em
$HHI = 1.298$. A contribuição americana domina o índice.

\medskip
\textbf{Jan--Mai 2026} (total exportado: US\$~703,7~mi):

\medskip
\begin{center}
\renewcommand{\arraystretch}{1.25}
\begin{tabular}{lS[table-format=3.1]S[table-format=2.2]S[table-format=3.1]}
\toprule
\textbf{Destino} & \textbf{Valor (US\$ mi)} & \textbf{Share (\%)} & \textbf{Share\textsuperscript{2}} \\
\midrule
Estados Unidos &  174.4 & 24.79 & 614.5 \\
México         &   53.6 &  7.62 &  58.1 \\
Argentina      &   44.8 &  6.37 &  40.6 \\
Itália         &   37.0 &  5.26 &  27.7 \\
Espanha        &   27.4 &  3.89 &  15.1 \\
Outros         &  366.5 & 52.07 & \textit{(vários)} \\
\midrule
\textbf{HHI total} & & & \textbf{840} \\
\bottomrule
\end{tabular}
\end{center}

\medskip
\noindent\textbf{O que explicou a queda?}
\[
  h_{\text{EUA}}^{2025} = 39{,}18^2 = 1535{,}1
  \quad\longrightarrow\quad
  h_{\text{EUA}}^{2026} = 24{,}79^2 = 614{,}5
  \quad\Rightarrow\quad
  \Delta h_{\text{EUA}} = -920{,}6 \text{ pontos}
\]
A simples queda do share americano de 39\% para 25\% respondeu por
\textbf{$\approx$~90\%} da variação total do índice ($\Delta HHI = -458$).
O crescimento de Itália, México e Espanha aumentou levemente o HHI
dos demais destinos, mas o efeito líquido foi a expressiva queda observada.
\end{exemplobox}

\subsection{Propriedades Importantes do HHI}

\begin{enumerate}
  \item \textbf{Sensibilidade à dominância:} o índice reage muito mais à
    redução de um destino grande do que ao crescimento de muitos destinos
    pequenos. Essa propriedade é exatamente o que ocorreu em 2026.

  \item \textbf{Independência do total exportado:} o HHI mede concentração
    nas \emph{participações}, não nos valores absolutos. Um setor pode
    exportar menos e ter HHI menor se sua pauta ficou mais diversificada.

  \item \textbf{Comparabilidade entre setores:} como utiliza participações
    percentuais, o HHI permite comparar setores de tamanhos muito diferentes
    (ex.: Papel e Celulose vs.\ Madeira).

  \item \textbf{Limitação --- qualidade dos destinos:} um HHI baixo pode
    refletir diversificação para mercados com baixo poder de compra ou alta
    instabilidade política. O índice deve ser complementado com análise
    qualitativa da pauta de destinos.
\end{enumerate}

%%=============================================================
\section{Metodologia 3 --- Estimativa de Receita Foregone: Análise Contrafactual}
\label{sec:receita}

\subsection{O que é uma Análise Contrafactual?}

Uma análise contrafactual responde à pergunta:
\emph{``o que teria acontecido se as tarifas não tivessem sido impostas?''}
Para isso, constrói-se um cenário hipotético --- o contrafactual ---
com base no comportamento histórico do setor \emph{antes} do choque tarifário.

A diferença entre o contrafactual (sem tarifas) e o observado (com tarifas)
é a \textbf{Receita Foregone} ($RF$) --- a receita que o setor
\emph{teria gerado} caso o choque não tivesse ocorrido.

\begin{interpretabox}[Receita Foregone $\neq$ Destruição Líquida]
\begin{itemize}[noitemsep]
  \item $DL_s$ compara o \emph{total exportado} nos dois períodos ---
    mede o que \textbf{de fato foi perdido} em valor total.
  \item $RF_s^k$ estima o que \textbf{teria sido exportado aos EUA}
    se a relação comercial fosse a do período $k$ ---
    mede a perda \emph{no canal americano}, mesmo que o total
    tenha crescido.
\end{itemize}
A $RF$ captura perda de participação no mercado americano mesmo em setores
cuja exportação total cresceu (como Papel e Celulose).
\end{interpretabox}

\subsection{Formulação Matemática}

A análise contrafactual é construída em duas etapas:

\subsubsection*{Etapa 1 --- Share de Referência}

Calcula-se o share histórico das exportações para os EUA no
período de referência $k$ (antes do choque):

\begin{formulabox}[Share de referência EUA]
\begin{equation}
  sh_{\text{EUA},s}^{k}
    \;=\; \frac{X_{\text{EUA},s}^{k}}{X_{\text{TOT},s}^{k}}
  \label{eq:share}
\end{equation}
\end{formulabox}

Esse share representa a proporção ``normal'' das exportações destinadas
aos EUA na ausência de tarifas.

\subsubsection*{Etapa 2 --- Exportações Contrafactuais e Receita Foregone}

Aplica-se o share de referência ao \emph{total efetivamente exportado
em 2026} (aceito como dado), obtendo a estimativa de quanto deveria
ter ido para os EUA:

\begin{formulabox}[Exportações contrafactuais e Receita Foregone]
\begin{align}
  \hat{X}_{\text{EUA},s}^{k}
    &\;=\; X_{\text{TOT},s}^{2026} \;\times\; sh_{\text{EUA},s}^{k}
    \label{eq:ctf} \\[6pt]
  RF_s^k
    &\;=\; \hat{X}_{\text{EUA},s}^{k} \;-\; X_{\text{EUA},s}^{2026}
    \label{eq:rf}
\end{align}
\end{formulabox}

\noindent\textbf{A intuição:} ``dado o total que o setor exportou em
Jan--Mai 2026, quanto disso deveria ter ido para os EUA, se a relação
comercial fosse a de antes?'' A diferença entre esse valor esperado e
o observado é a receita perdida no mercado americano.

\subsection{Os Três Cenários de Referência}

A análise utiliza três horizontes históricos, cada um respondendo a uma
pergunta diferente:

\begin{center}
\renewcommand{\arraystretch}{1.4}
\begin{tabular}{>{\bfseries}clp{7.5cm}}
\toprule
Cenário & Período $k$ & Pergunta que responde \\
\midrule
A & Jan--Mai 2025
  & Se nada mudasse em relação ao primeiro semestre de 2025, quanto
    iria aos EUA? \emph{(conservador --- já embute antecipação às tarifas)} \\[0.3em]
B & 2024 (ano completo)
  & Se a relação de 2024 --- o ano imediatamente anterior ao choque ---
    fosse mantida, qual seria a perda? \emph{(referência principal)} \\[0.3em]
C & 2023 (ano completo)
  & Se a dependência estrutural pré-Trump fosse a referência, qual a
    perda máxima? \emph{(limite superior)} \\
\bottomrule
\end{tabular}
\end{center}

\subsubsection*{Por que três cenários?}

Nenhum cenário isolado é perfeito:

\begin{itemize}
  \item O Cenário~A (Jan-Mai 2025) já pode refletir antecipação das tarifas
    por parte de importadores e exportadores. Tende a \textbf{subestimar} a perda.

  \item O Cenário~B (2024 completo) é o mais neutro: representa o período
    imediatamente pré-choque, sem antecipação significativa. É o
    \textbf{cenário de referência principal}.

  \item O Cenário~C (2023) captura a relação histórica antes do retorno
    das tensões comerciais, servindo como \textbf{teto máximo} da estimativa.
\end{itemize}

\noindent A \textbf{convergência entre B e C} (US\$~114,4 vs.\ US\$~114,5~mi
para o Complexo Florestal) é um sinal de \textbf{robustez metodológica}:
indica que a estrutura de dependência dos EUA era estável em 2023--2024,
e que o intervalo estimado é confiável.

\subsection{Exemplo Passo a Passo: Complexo Florestal (Cenário B)}

\begin{exemplobox}[Receita Foregone --- Complexo Florestal, Cenário B (base 2024)]
\small
\textbf{Dados de entrada:}
\begin{align*}
  X_{\text{EUA},\text{CF}}^{2024\text{ (ano)} } &= \text{US\$~}1{.}363{,}7 \text{ mi}
  &\quad X_{\text{TOT},\text{CF}}^{2024\text{ (ano)} } &= \text{US\$~}3{.}321{,}4 \text{ mi} \\
  X_{\text{EUA},\text{CF}}^{2026\text{ (jan-mai)}} &= \text{US\$~}174{,}4 \text{ mi}
  &\quad X_{\text{TOT},\text{CF}}^{2026\text{ (jan-mai)}} &= \text{US\$~}703{,}7 \text{ mi}
\end{align*}

\textbf{Passo 1 --- Share de referência (2024):}
\[
  sh_{\text{EUA},\text{CF}}^{B}
    = \frac{1{.}363{,}7}{3{.}321{,}4}
    = 41{,}04\%
\]

\textbf{Passo 2 --- Exportações contrafactuais:}
\[
  \hat{X}_{\text{EUA},\text{CF}}^{B}
    = 703{,}7 \times 0{,}4104
    = \text{US\$~} 288{,}8 \text{ mi}
\]

\textbf{Passo 3 --- Receita Foregone:}
\[
  RF_{\text{CF}}^{B}
    = 288{,}8 - 174{,}4
    = \mathbf{\text{US\$~} 114{,}4 \text{ milhões}}
\]

\medskip
\noindent\textit{Interpretação:} Se o Complexo Florestal tivesse mantido
em Jan--Mai 2026 o mesmo share dos EUA de 2024 (41\%), teria exportado
US\$~288,8~mi para aquele mercado.
Como exportou apenas US\$~174,4~mi, a receita foregone é de
US\$~114,4~mi --- equivalente a
\textbf{16,3\% do total exportado no período}.
\end{exemplobox}

%%=============================================================
\section{Resultados por CNAE --- Jan-Mai 2025 vs.\ 2026}
\label{sec:resultados}

\subsection{Decomposição do Impacto Tarifário}

\begin{table}[H]
  \centering
  \caption{Decomposição do impacto tarifário por CNAE (US\$ milhões)}
  \label{tab:decomp}
  \renewcommand{\arraystretch}{1.4}
  \small
  \begin{tabular}{l
      S[table-format=3.1]
      S[table-format=2.1]
      S[table-format=3.1]
      S[table-format=2.1]
      S[table-format=3.1]}
    \toprule
    \textbf{Setor}
      & \textbf{$DB_s$}
      & \textbf{$DC_s$}
      & \textbf{$DL_s$}
      & \textbf{$EP_s$}
      & \textbf{$PT_s$} \\
    \midrule
    Madeira (CNAE 16)          & 118.1 & 32.5 &  85.6 & 30.1 & 115.7 \\
    Móveis (CNAE 31)           &  24.9 &  0.0 &  24.9 &  0.0 &  24.9 \\
    Papel e Celulose (CNAE 17) &   5.2 & 10.8 &  -5.6 & 17.7 &  12.2 \\
    Base Florestal (CNAE 2)    &  -0.6 &  0.0 &  -0.6 &  0.0 &  -0.6 \\
    \midrule
    \textbf{Complexo Florestal}
      & \textbf{147.7} & \textbf{29.2} & \textbf{118.6}
      & \textbf{0.0}   & \textbf{118.6} \\
    \bottomrule
  \end{tabular}
  \begin{minipage}{\textwidth}
    \vspace{4pt}\footnotesize
    \textit{DB} = Destruição Bruta \;|\;
    \textit{DC} = Desvio de Comércio \;|\;
    \textit{DL} = Destruição Líquida \;|\;
    \textit{EP} = Efeito-Preço \;|\;
    \textit{PT} = Perda Total.\\
    O $EP$ do Complexo consolidado é zero porque o volume total para outros mercados
    \emph{recuou} no período, apesar do valor em US\$ ter crescido
    (evidência de aumento de preço médio sem expansão de volume).
    Valores negativos em $DB$ (Base Florestal) indicam crescimento
    das exportações para os EUA no período.
  \end{minipage}
\end{table}

\subsection{Evolução do HHI de Concentração de Destinos}

\begin{table}[H]
  \centering
  \caption{Evolução do HHI por CNAE --- Jan-Mai 2024 a Jan-Mai 2026}
  \label{tab:hhi}
  \renewcommand{\arraystretch}{1.4}
  \small
  \begin{tabular}{l
      S[table-format=4.0]
      S[table-format=4.0]
      S[table-format=4.0]
      S[table-format=-4.0]
      l}
    \toprule
    \textbf{Setor}
      & \textbf{HHI 2024}
      & \textbf{HHI 2025}
      & \textbf{HHI jan-mai 26}
      & \textbf{$\Delta$ 24--26}
      & \textbf{Classificação 2026} \\
    \midrule
    Madeira (CNAE 16)          & 2695 & 1929 & 1253 & -1442 & Baixa concentração \\
    Móveis (CNAE 31)           & 2362 & 1877 & 1314 & -1048 & Baixa concentração \\
    Papel e Celulose (CNAE 17) & 1059 & 1118 & 1078 &   +19 & Baixa concentração \\
    \midrule
    \textbf{Complexo Florestal}
      & \textbf{1863} & \textbf{1298} & \textbf{840}
      & \textbf{-1023} & \textbf{Diversificado} \\
    \bottomrule
  \end{tabular}
  \begin{minipage}{\textwidth}
    \vspace{4pt}\footnotesize
    HHI calculado sobre valor FOB de Jan--Mai de cada ano (base comparável).
    Papel e Celulose manteve HHI estável (+19 pontos) pois não sofreu
    tarifa de 50\% --- sua concentração em 2024 já era baixa (1.059).
    $\Delta = $ HHI jan-mai 2026 $-$ HHI 2024.
  \end{minipage}
\end{table}

\subsection{Estimativa da Receita Foregone por Cenário}

\begin{table}[H]
  \centering
  \caption{Receita Foregone por cenário e CNAE --- Jan-Mai 2026 (US\$ milhões)}
  \label{tab:ctf}
  \renewcommand{\arraystretch}{1.4}
  \small
  \begin{tabular}{l
      S[table-format=2.2]S[table-format=3.1]
      S[table-format=2.2]S[table-format=3.1]
      S[table-format=2.2]S[table-format=3.1]}
    \toprule
    \multirow{2}{*}{\textbf{Setor}}
      & \multicolumn{2}{c}{\textbf{Cenário A} (jan-mai 2025)}
      & \multicolumn{2}{c}{\textbf{Cenário B} (2024 completo)}
      & \multicolumn{2}{c}{\textbf{Cenário C} (2023 completo)} \\
    \cmidrule(lr){2-3}\cmidrule(lr){4-5}\cmidrule(lr){6-7}
      & \textbf{$sh$ EUA (\%)} & \textbf{RF (mi)}
      & \textbf{$sh$ EUA (\%)} & \textbf{RF (mi)}
      & \textbf{$sh$ EUA (\%)} & \textbf{RF (mi)} \\
    \midrule
    Madeira (CNAE 16)          & 48.58 &  76.5 & 50.11 &  83.5 & 50.65 &  85.9 \\
    Móveis (CNAE 31)           & 43.35 &  10.5 & 45.68 &  12.3 & 47.64 &  13.9 \\
    Papel e Celulose (CNAE 17) &  7.72 &   5.6 &  7.61 &   5.5 & 10.35 &   9.7 \\
    Base Florestal (CNAE 2)    &  0.47 &  -0.6 &  3.63 &  -0.1 &  3.81 &   0.0 \\
    \midrule
    \textbf{Complexo Florestal}
      & \textbf{39.18} & \textbf{101.3}
      & \textbf{41.04} & \textbf{114.4}
      & \textbf{41.06} & \textbf{114.5} \\
    \bottomrule
  \end{tabular}
  \begin{minipage}{\textwidth}
    \vspace{4pt}\footnotesize
    \textit{sh EUA} = share de referência para os EUA no período $k$.
    \textit{RF} = Receita Foregone (US\$ mi, arredondado).
    Valores negativos (Base Florestal, CNAE~2) indicam que as exportações para
    os EUA \emph{superaram} o contrafactual no período --- o setor ganhou
    participação americana.
    A convergência entre Cenários B e C para o Complexo Florestal
    (US\$~114,4 vs.\ US\$~114,5~mi) valida a robustez do intervalo estimado.
  \end{minipage}
\end{table}

%%=============================================================
\section{Síntese e Limitações Metodológicas}

\subsection{Síntese por Setor}

\begin{description}[style=nextline]
  \item[Madeira (CNAE 16) --- impacto severo, parcialmente mitigado]
    Maior destruição bruta do Complexo (US\$~118,1~mi). O setor encontrou
    mercados alternativos (+US\$~32,5~mi de desvio), mas o alto prêmio EUA
    (US\$~0,51/kg; +117\%) gerou efeito-preço de US\$~30,1~mi.
    Perda total: \textbf{US\$~115,7~mi}.
    HHI caiu de 2.695 para 1.253 --- a maior redução de concentração
    observada no Complexo.

  \item[Móveis (CNAE 31) --- sem desvio, vulnerabilidade estrutural]
    Perda de US\$~24,9~mi sem nenhuma compensação nos demais mercados
    ($DC = EP = 0$). A ausência de desvio indica que os mercados alternativos
    não absorveram volume equivalente ao perdido nos EUA.
    HHI ainda em concentração moderada (1.314) em 2026.

  \item[Papel e Celulose (CNAE 17) --- resiliência tarifária]
    Único setor com tarifa parcial (10\% vs.\ 50\%). A queda nos EUA
    (US\$~5,2~mi) foi superada pelo crescimento em outros destinos
    (+US\$~10,8~mi), resultando em destruição líquida negativa.
    Contudo, o efeito-preço (US\$~17,7~mi) revela que o crescimento se
    deu a preços inferiores. HHI estável.

  \item[Base Florestal (CNAE 2) --- impacto desprezível]
    Exportações para os EUA cresceram ($DB = -$US\$~0,6~mi),
    possivelmente por contratos pré-tarifas.

  \item[Complexo Florestal (consolidado)]
    Destruição bruta de US\$~147,7~mi; desvio de US\$~29,2~mi;
    destruição líquida de US\$~118,6~mi.
    Receita foregone entre \textbf{US\$~101~mi} (conservador)
    e \textbf{US\$~115~mi} (base/estrutural).
    HHI passou de 1.863 (2024) para 840 (Jan-Mai 2026) ---
    de concentração moderada para diversificado em 18 meses.
\end{description}

\subsection{Limitações Metodológicas}

\begin{enumerate}
  \item \textbf{Causalidade vs.\ correlação:} a metodologia atribui toda a
    queda nas exportações para os EUA às tarifas. Outros fatores (ciclo da
    construção civil americana, câmbio, estoques de importadores) podem ter
    contribuído.

  \item \textbf{Desvio medido em valor, não em quantidade:} o ganho em
    mercados alternativos pode refletir aumento de preço, não de volume.
    O Efeito-Preço tenta corrigir isso, mas assume prêmio uniforme
    por CNAE (calculado com mix fixo de produtos).

  \item \textbf{Share de referência estático:} o contrafactual supõe que
    a tendência de share dos EUA seria constante. Ignora possível crescimento
    orgânico que ocorreria no mercado americano em 2026 sem tarifas.

  \item \textbf{Agregação por CNAE:} a heterogeneidade dentro de um setor
    (ex.: madeira serrada vs.\ MDF dentro de CNAE~16) pode gerar resultados
    distintos ao nível de produto. A análise ao nível NCM-8 pode ser
    desenvolvida como extensão.

  \item \textbf{Período de análise:} cinco meses (Jan-Mai 2026) capturam o
    impacto imediato. Efeitos de médio prazo --- fechamento de empresas,
    perda de contratos de longo prazo, realocação de capacidade produtiva
    --- não estão refletidos.

  \item \textbf{HHI e qualidade dos destinos:} o índice mede dispersão
    geográfica, não qualidade comercial dos novos destinos. Um HHI baixo
    pode mascarar dependência de mercados com menor poder de compra
    ou maior instabilidade.
\end{enumerate}

%%=============================================================
\section*{Referências}
\addcontentsline{toc}{section}{Referências}

\begin{enumerate}[label={[\arabic*]}]
  \item BRASIL. Ministério do Desenvolvimento, Indústria, Comércio e
    Serviços (MDIC). \textit{ComexStat --- Sistema de Análise das Informações
    de Comércio Exterior}. Brasília, 2026.
    Disponível em: \url{https://comexstat.mdic.gov.br}.

  \item FIESC. \textit{Dicionário NCM--CNAE --- Mapeamento do Complexo
    Florestal de Santa Catarina}. Versão atualizada em 07/04/2026.
    Florianópolis: Observatório FIESC, 2026.

  \item HERFINDAHL, Orris C. \textit{Concentration in the Steel Industry}.
    Dissertação (Doutorado) --- Columbia University, 1950.

  \item HIRSCHMAN, Albert O. \textit{National Power and the Structure of
    Foreign Trade}. Berkeley: University of California Press, 1945.

  \item U.S.\ Department of Justice; Federal Trade Commission.
    \textit{Horizontal Merger Guidelines}. Washington, DC: DOJ/FTC, 2010.
    Critérios de classificação do HHI adotados na Seção~\ref{sec:hhi}.

  \item VINER, Jacob. \textit{The Customs Union Issue}. New York: Carnegie
    Endowment for International Peace, 1950.
    Referência conceitual para \textit{trade creation} e
    \textit{trade diversion}.

  \item WORLD BANK. \textit{Export Diversification and Quality}.
    Washington, DC: World Bank, 2012.
    Uso do HHI para mensurar diversificação de destinos de exportação.
\end{enumerate}

%%=============================================================
\appendix

\section{Glossário Completo de Variáveis}
\label{app:glossario}

\begin{center}
\renewcommand{\arraystretch}{1.4}
\begin{tabular}{>{\bfseries}lp{10.5cm}}
\toprule
Símbolo & Descrição \\
\midrule
$X_{d,s}^{t}$
  & Valor FOB (US\$) exportado pelo setor $s$ para o destino $d$
    em Jan--Mai do ano $t$ \\
$Q_{d,s}^{t}$
  & Volume (kg líquido) exportado pelo setor $s$ para $d$
    em Jan--Mai do ano $t$ \\
$X_{\text{TOT},s}^{t}$
  & Total exportado pelo setor $s$ em Jan--Mai do ano $t$
    (soma sobre todos os destinos) \\
$VU_{d,s}^{t}$
  & Valor unitário médio: $X_{d,s}^{t}/Q_{d,s}^{t}$ (US\$/kg) \\
EUA
  & Estados Unidos (código ISO 249) \\
OUT
  & Todos os demais países (complemento de EUA) \\
CF
  & Complexo Florestal (CNAE 2 + 16 + 17 + 31) \\
$DB_s$
  & Destruição Bruta: $X_{\text{EUA},s}^{2025} - X_{\text{EUA},s}^{2026}$ \\
$DC_s$
  & Desvio de Comércio:
    $\max\bigl(X_{\text{OUT},s}^{2026} - X_{\text{OUT},s}^{2025},\;0\bigr)$ \\
$DL_s$
  & Destruição Líquida: $DB_s - DC_s$ \\
$\pi_s$
  & Prêmio de Mercado EUA:
    $VU_{\text{EUA},s}^{2025} - VU_{\text{OUT},s}^{2025}$ (US\$/kg) \\
$\Delta Q_{\text{OUT},s}$
  & Variação de volume para OUT:
    $Q_{\text{OUT},s}^{2026} - Q_{\text{OUT},s}^{2025}$ \\
$EP_s$
  & Efeito-Preço:
    $\max(\Delta Q_{\text{OUT},s},\,0)\cdot\max(\pi_s,\,0)$ \\
$PT_s$
  & Perda Total: $DL_s + EP_s$ \\
$HHI_{s,t}$
  & Índice Herfindahl-Hirschman do setor $s$ no período $t$ \\
$h_{i,s}^{t}$
  & Contribuição individual do país $i$ ao HHI:
    $\bigl(X_{i,s}^{t}/X_{\text{TOT},s}^{t}\times 100\bigr)^2$ \\
$N_{s,t}$
  & Número de países destino com exportação positiva
    (setor $s$, período $t$) \\
$sh_{\text{EUA},s}^{k}$
  & Share EUA de referência no cenário contrafactual $k$:
    $X_{\text{EUA},s}^{k}/X_{\text{TOT},s}^{k}$ \\
$\hat{X}_{\text{EUA},s}^{k}$
  & Exportações contrafactuais para os EUA no cenário $k$:
    $X_{\text{TOT},s}^{2026} \times sh_{\text{EUA},s}^{k}$ \\
$RF_s^k$
  & Receita Foregone no cenário $k$:
    $\hat{X}_{\text{EUA},s}^{k} - X_{\text{EUA},s}^{2026}$ \\
\bottomrule
\end{tabular}
\end{center}

\end{document}
